\chapter{Tooth extraction}
\index{Tooth extraction}

\section*{Definition and Overview}
Tooth extraction, clinically termed exodontia, refers to the painless surgical removal of a tooth or root remnants from its bony socket (alveolus) within the alveolar process of either the maxilla (upper jaw) or mandible (lower jaw). While contemporary dentistry emphasizes the conservation of natural dentition through advanced endodontic, periodontic, and restorative techniques, tooth extraction remains one of the most fundamental, ancient, and frequently performed procedures in oral surgery.

The indications for extraction are diverse, ranging from advanced dental caries that have compromised the tooth structure beyond restorability, to severe periodontal disease causing irreversible loss of supporting bone, dental trauma, orthodontic spatial requirements, and the management of impacted teeth (particularly third molars).

Extractions are broadly classified into two categories based on surgical complexity: simple (or closed) extractions and surgical (or open) extractions. A simple extraction is performed on a tooth that is visible in the oral cavity, typically using dental elevators and forceps under local anaesthesia to luxate and lift the tooth from its socket. Conversely, a surgical extraction is required for teeth that cannot be easily accessed, such as those fractured at the gumline, severely decayed, or impacted within the bone. This procedure involves the elevation of a mucoperiosteal flap, bone removal (ostectomy), and frequently, sectioning of the tooth (odontosection) to facilitate removal.

Modern dental practice views extraction not merely as an end-point, but as the initial stage of oral rehabilitation. Consequently, maintaining the integrity of the surrounding alveolar bone through meticulous surgical technique and socket preservation has become integral to the procedure, preparing the site for subsequent rehabilitation with dental implants, fixed bridges, or removable prostheses.

\section*{Epidemiology}
Tooth extraction is a ubiquitous procedure with significant epidemiological variation across different demographics, age groups, and socioeconomic strata. Globally, dental caries and periodontal disease remain the primary drivers of tooth loss, leading to hundreds of millions of extractions annually.

The distribution of extractions by age exhibits distinct peaks. In adolescents and young adults (typically aged 17--25 years), the extraction of third molars (wisdom teeth) is the most common surgical intervention. Studies indicate that up to 60\% to 70\% of individuals in developed countries undergo extraction of at least one third molar due to impaction, pericoronitis, or orthodontic requirements. Conversely, in mature and geriatric populations, tooth extraction is primarily necessitated by periodontal disease and chronic, untreated dental caries. By the age of 65, approximately 60\% of the global population has lost at least one permanent tooth (excluding third molars), and severe tooth loss (having fewer than 9 remaining teeth) affects approximately 5\% to 10\% of the elderly population.

Socioeconomic status is a major determinant of extraction prevalence. Individuals in lower socioeconomic brackets or residing in regions with limited access to fluoridated water and dental care exhibit higher rates of extraction. In these cohorts, dental pain is more frequently managed via emergency extraction rather than expensive, tooth-saving restorative procedures such as root canal therapy, crowns, or periodontal surgery. Furthermore, geographic variations show that in low-income countries, extraction is often the only available dental treatment, whereas in high-income nations, the focus has shifted towards preventative care, resulting in a gradual decline in the overall rate of extractions per capita.

\section*{Aetiology and Pathophysiology}
The necessity for tooth extraction arises from various pathological processes that compromise the structural, pulpal, or periodontal integrity of the tooth, or from systemic conditions that alter the dental arch. The primary etiologies and their underlying pathophysiological mechanisms include:

\begin{itemize}
  \item **Severe Dental Caries (Tooth Decay):** Chronic bacterial infection primarily initiated by \textit{Streptococcus mutans} and \textit{Lactobacillus} species. These micro-organisms ferment dietary carbohydrates, producing organic acids that cause the demineralization of enamel and dentine. As the carious lesion progresses, it penetrates the pulp chamber, leading to microvascular congestion, irreversible pulpitis, pulpal necrosis, and subsequent periapical pathology (such as periapical granulomas or abscesses). If the remaining coronal tooth structure is insufficient to support a restoration or crown, extraction is indicated.
  \item **Advanced Periodontitis:** A chronic inflammatory disease initiated by subgingival plaque biofilms containing periodontal pathogens (such as \textit{Porphyromonas gingivalis}). The host's immunoinflammatory response leads to the progressive destruction of the periodontal ligament (PDL), cementum, and alveolar bone. The loss of these supporting tissues leads to clinical attachment loss, pocket formation, and pathological tooth mobility (Class III mobility), rendering the tooth non-functional and susceptible to recurrent periodontal abscesses.
  \item **Pathological Tooth Impaction:** The failure of a tooth to erupt fully into its functional position in the dental arch within the expected developmental timeline. Impaction is most common in mandibular and maxillary third molars, followed by maxillary canines. The primary causes include insufficient dental arch length, abnormal eruption pathways, or physical obstruction by adjacent teeth. Pathophysiologically, partially impacted teeth create a plaque-retentive area under the overlying operculum, leading to pericoronitis, localized caries on the distal aspect of the adjacent second molar, external root resorption, or the development of odontogenic cysts (e.g., dentigerous cysts) from the dental follicle.
  \item **Dentoalveolar Trauma:** High-energy impact injuries (e.g., sports injuries, falls, motor vehicle accidents) that cause tooth fracture, luxation, or intrusion. Vertical root fractures, which extend longitudinally from the crown down the root surface, have a hopeless prognosis. These fractures create a direct pathway for oral bacteria to colonize the periodontal ligament space, causing rapid localized osteitis, bone loss, and chronic infection, necessitating extraction.
  \item **Failed Endodontic Treatment:** Persistent or recurrent pulpal/periapical pathology after root canal therapy. This occurs due to complex root canal anatomy (such as accessory canals or severe curvatures) that prevents adequate debridement, procedural errors (e.g., separated instruments, root perforations), or microleakage of the coronal restoration, allowing re-infection of the canal space. When endodontic retreatment or surgical apicoectomy is impossible or fails, extraction is required.
  \item **Orthodontic Spatial Requirements:** Extraction of permanent teeth (most commonly first or second premolars) to resolve severe dental crowding, facilitate the retraction of protrusive anterior teeth, or correct skeletal discrepancies (such as Class II or Class III malocclusions).
  \item **Ectopic and Supernumerary Teeth:** Teeth that erupt in abnormal positions or are present in excess of the normal dental formula (e.g., mesiodens). These teeth can prevent the eruption of normal succedaneous teeth, cause root resorption of adjacent teeth, or cause functional and occlusal interference.
  \item **Systemic and Therapeutic Risk Factors:** Patients undergoing head and neck radiation therapy are at risk of osteoradionecrosis (ORN) if teeth within the radiation field are extracted post-therapy; thus, pre-radiation extractions of teeth with poor prognosis are mandatory. Similarly, patients taking antiresorptive medications (e.g., bisphosphonates, denosumab) are at risk of Medication-Related Osteonecrosis of the Jaw (MRONJ), requiring careful pre-operative planning and conservative extraction techniques.
\end{itemize}

\section*{Clinical Features}
The clinical presentation of a patient requiring tooth extraction varies depending on the underlying pathology. A detailed evaluation of these clinical features is essential for formulating a correct diagnosis and planning the extraction procedure:

\begin{itemize}
  \item **Odontogenic Pain:** The most common symptom. It can present as sharp, transient pain triggered by thermal, chemical, or osmotic stimuli (typical of reversible pulpitis), or as a continuous, deep, throbbing pain that is spontaneous, poorly localized, and radiates to the ear, temple, or neck (indicative of irreversible pulpitis or acute periapical infection). Pain upon percussion or mastication indicates that the inflammatory process has spread beyond the pulp chamber into the periodontal ligament.
  \item **Pathological Tooth Mobility:** Marked movement of the tooth within its socket. Mobility is assessed clinically by applying horizontal and vertical force using the ends of two instruments. It is graded from Class I (horizontal movement up to 1~mm), Class II (horizontal movement exceeding 1~mm), to Class III (severe horizontal mobility exceeding 1~mm accompanied by vertical depressibility within the socket).
  \item **Gingival and Mucosal Swelling:** Erythema, swelling, and tenderness of the soft tissues surrounding the affected tooth. In pericoronitis, the soft tissue flap (operculum) covering the partially erupted third molar becomes swollen, red, and tender, occasionally displaying imprints of the opposing maxillary tooth.
  \item **Suppuration (Pus Discharge):** Active purulent drainage from the gingival sulcus or the presence of a mucosal sinus tract (fistula) with a draining parulis (gum boil). This indicates a chronic periapical or periodontal abscess that has breached the cortical bone.
  \item **Trismus:** Limited mouth opening (typically a maximum interincisal distance of less than 35~mm) caused by inflammatory spasm of the masticatory muscles. This is common when infections from mandibular molars spread into the masseteric, pterygomandibular, or submandibular fascial spaces.
  \item **Halitosis and Dysgeusia:** A foul, offensive breath odor and a persistent bad, metallic, or bitter taste in the mouth. This is caused by the proliferation of anaerobic bacteria, necrotic pulpal tissue, and the accumulation of food debris in carious cavities, periodontal pockets, or under an operculum.
  \item **Systemic Signs of Spreading Infection:** Generalized pyrexia (fever), rigors, malaise, headache, and regional lymphadenopathy (swollen, tender submandibular or upper deep cervical lymph nodes), indicating that the localized dental infection is spreading systemically.
\end{itemize}

\section*{Differential Diagnosis}
Before proceeding with an irreversible procedure like tooth extraction, clinicians must differentiate odontogenic pain requiring extraction from other dental or non-dental conditions that can be managed conservatively:

\begin{itemize}
  \item **Reversible Pulpitis:** Mild to moderate inflammation of the dental pulp. The pain is sharp and transient, lasting only seconds after the removal of a stimulus (cold, hot, or sweet). The pulp remains vital, and the condition is fully reversible through restorative treatments (such as dental fillings or direct pulp capping), preserving the tooth.
  \item **Dentinal Hypersensitivity:** Sharp, transient pain arising from exposed dentinal tubules, commonly due to gingival recession or enamel erosion. The pain is triggered by external stimuli (thermal, chemical, or tactile) and is managed non-invasively with desensitizing toothpastes, dentine bonding agents, or gingival grafting.
  \item **Trigeminal Neuralgia:** A neuropathic disorder characterized by sudden, severe, unilateral, stabbing, or shock-like pain along the distribution of the trigeminal nerve. The pain is brief, lasting from a few seconds to two minutes, and is triggered by light touch, washing, shaving, or brushing teeth, occurring in the absence of clinical or radiographic dental pathology.
  \item **Temporomandibular Disorders (TMD):** A group of musculoskeletal conditions affecting the temporomandibular joint (TMJ) and the muscles of mastication. Symptoms include dull, aching pain in the preauricular region, joint clicking or crepitus, and restricted jaw movement, which can mimic lower molar pain but is unrelated to dental pathology.
  \item **Acute Maxillary Sinusitis:** Inflammation of the maxillary sinus mucosa, often secondary to an upper respiratory tract infection. Because the roots of the maxillary premolars and molars are in close proximity to the sinus floor, sinus pressure can refer pain to multiple upper teeth. This pain is typically bilateral, exacerbated by bending forward or walking, and accompanied by nasal congestion and pressure.
  \item **Atypical Odontalgia (Phantom Tooth Pain):** A chronic neuropathic pain condition characterized by continuous, daily throbbing pain in a tooth or an extraction site that has no identifiable clinical or radiographic cause. Invasive dental treatments, including root canal therapy or extraction, fail to relieve the pain and often exacerbate it.
  \item **Myofascial Pain Referral:** Pain referred to the teeth from active trigger points in the masticatory muscles, particularly the masseter and temporalis muscles. Palpation of the muscle trigger points reproduces the dental pain, which can be managed with physical therapy, trigger point injections, or occlusal splints.
\end{itemize}

\section*{When to Seek Medical Care}
Although tooth extraction is typically an elective, outpatient procedure, dental infections can occasionally progress rapidly, leading to life-threatening complications such as airway obstruction or systemic sepsis. Patients must be educated on the warning signs that necessitate immediate medical evaluation.

Emergency services (such as local emergency services in many countries, 911 in the United States, or 999 in the United Kingdom) should be contacted immediately if any of the following symptoms occur:
\begin{itemize}
  \item **Difficulty Breathing (Dyspnoea) or Swallowing (Dysphagia):** Indicates rapid swelling of the submandibular, sublingual, or lateral pharyngeal spaces (e.g., Ludwig's angina), which can compromise the airway.
  \item **Inability to Control Saliva:** Drooling indicates an inability to swallow, which is a sign of impending airway obstruction.
  \item **Rapidly Spreading Facial or Neck Swelling:** Particularly if the swelling is firm, painful, warm, and extends to the cheek, eyelid, or down the neck.
  \item **High Fever (above 38.5 degrees Celsius or 101.3 degrees Fahrenheit):** Especially when accompanied by chills, rigors, confusion, severe lethargy, or a rapid heart rate, suggesting systemic bacteraemia or sepsis.
  \item **Uncontrolled Bleeding:** Continuous, profuse bleeding from the extraction site that does not slow down or stop after applying firm, direct pressure with a clean gauze pad for 30 to 60 minutes.
  \item **Severe, Unrelenting Pain:** Pain that is uncontrolled by maximum doses of prescribed analgesics, which may indicate a severe infection (osteomyelitis) or dry socket.
\end{itemize}

\section*{Diagnosis and Investigations}
A thorough diagnostic workup is required before any extraction to evaluate the tooth, its surrounding structures, and the patient's general health:

\begin{itemize}
  \item **Clinical Examination:** Includes visual inspection of the tooth and surrounding soft tissues, assessment of dental caries, tooth mobility, gingival inflammation, pericoronal tissue health, and the presence of mucosal swellings or fistulous tracts.
  \item **Periodontal Probing and Mobility Assessment:** Measurement of periodontal pocket depths using a periodontal probe. Depths greater than 5~mm indicate attachment loss, while depths exceeding 8~mm associated with Class III mobility confirm terminal periodontal disease with a hopeless prognosis.
  \item **Percussion and Palpation Testing:** Tapping the tooth vertically and horizontally with an instrument handle to assess periodontal ligament inflammation. Palpating the apical area of the tooth checks for cortical plate expansion or fluctuant swelling.
  \item **Pulp Vitality Testing:** Thermal tests (using cold spray or hot gutta-percha) and electric pulp testing (EPT) are performed to determine the sensory response of the pulp. A lack of response suggests pulp necrosis.
  \item **Intraoral Periapical Radiography (IOPA):** The primary imaging modality. It provides high-resolution details of the tooth crown, pulp chamber, root canal system, root number, shape, and curvature, periodontal ligament space, and surrounding alveolar bone. It reveals the extent of bone loss, periapical radiolucencies (abscesses, cysts, or granulomas), and root fractures.
  \item **Orthopantomography (OPG):** A panoramic radiograph showing both jaws, TMJs, and sinuses. It is essential for planning wisdom teeth extractions, evaluating generalized periodontal disease, and identifying the relationship between the roots of mandibular molars and the inferior alveolar nerve canal.
  \item **Cone Beam Computed Tomography (CBCT):** Three-dimensional imaging used for complex cases. CBCT is indicated when a panoramic radiograph shows close proximity or overlap between the roots of a mandibular third molar and the inferior alveolar nerve canal, allowing the surgeon to assess the relationship in three dimensions to prevent nerve damage.
  \item **Coagulation Profiling:** Laboratory tests including Prothrombin Time (PT), International Normalized Ratio (INR), and Activated Partial Thromboplastin Time (aPTT) for patients on oral anticoagulants (e.g., Warfarin) or with suspected coagulopathy. An INR of 2.0--3.0 is typically acceptable for simple extractions, but consultation with the patient's physician is mandatory.
  \item **Complete Blood Count (CBC):** Indicated in patients with systemic symptoms (fever, malaise) to assess white blood cell counts and evaluate the severity of active infection.
\end{itemize}

\section*{Management}
The management of a tooth extraction encompasses pre-operative preparation, the surgical procedure itself (simple or surgical), and post-operative care:

\begin{itemize}
  \item **Pre-Operative Preparation:** Involves reviewing the patient's medical history (focusing on cardiovascular disease, bleeding disorders, medications, and allergies), obtaining informed consent, and discussing the risks, benefits, and alternatives. Patients at risk for infective endocarditis (e.g., those with prosthetic heart valves or a history of endocarditis) are prescribed prophylactic antibiotics (e.g., Amoxicillin 2~g or Clindamycin 600~mg) 1 hour before the procedure.
  \item **Anaesthesia and Pain Management:** Local anaesthesia is standard, using agents like Lidocaine (2\% with 1:100,000 epinephrine) or Articaine (4\% with 1:100,000 epinephrine). Maxillary teeth are anaesthetized via infiltration, while mandibular teeth require nerve blocks (e.g., inferior alveolar nerve block). Nitrous oxide inhalation sedation, intravenous (IV) conscious sedation, or general anaesthesia may be used for anxious patients or complex surgical extractions.
  \item **Simple (Closed) Extraction Technique:** Executed in four stages: (1) Severing the epithelial attachment around the tooth neck using a periosteal elevator. (2) Luxating the tooth by inserting elevators (e.g., straight elevators) into the periodontal ligament space, applying slow, controlled rotational and leverage forces to expand the alveolar bone socket and tear the PDL. (3) Applying extraction forceps to the tooth root and executing apical, buccal, and lingual/palatal movements to further expand the socket. (4) Delivery of the tooth from the socket.
  \item **Surgical (Open) Extraction Technique:** Required for impacted, severely decayed, or fractured teeth. It involves: (1) Reflecting a mucoperiosteal flap using a scalpel and periosteal elevator. (2) Conserving bone removal (ostectomy) under sterile saline irrigation to prevent bone necrosis. (3) Tooth sectioning (odontosection) into individual roots to facilitate removal. (4) Socket curettage to remove granulation tissue or cyst fragments, smoothing sharp bone edges, saline irrigation, flap repositioning, and suturing with resorbable sutures.
  \item **Socket Preservation (Alveolar Ridge Preservation):** Placing a bone graft material (allograft, xenograft, or synthetic bone substitute) into the fresh extraction socket, covered with a collagen membrane. This minimizes bone resorption and maintains bone volume for future dental implant placement.
  \item **Post-Operative Pharmacotherapy:** Non-opioid analgesics are the gold standard: NSAIDs (Ibuprofen 400--600~mg) combined with Paracetamol (500--1000~mg) every 6 hours. Routine post-operative antibiotics are contraindicated in healthy patients.
\end{itemize}

\section*{Complications}
Tooth extraction is a surgical procedure associated with several potential intraoperative and post-operative complications:

\begin{itemize}
  \item **Alveolar Osteitis (Dry Socket):** The most common post-operative complication, occurring in 1\% to 5\% of all extractions and up to 30\% of mandibular third molar extractions. It occurs when the blood clot in the socket prematurely disintegrates or fails to form, exposing the underlying alveolar bone to oral bacteria and debris. This causes intense, throbbing pain radiating to the ear, typically starting 3--4 days post-extraction, accompanied by halitosis. Risk factors include smoking, oral contraceptive use, poor oral hygiene, and surgical trauma.
  \item **Post-Operative Haemorrhage:** Excessive bleeding from the socket. Primary (intraoperative), reactionary (2--6 hours post-op as epinephrine vasoconstriction fades), or secondary (1--2 weeks post-op due to infection). Managed with gauze pressure, gelatin sponges, tranexamic acid mouthwashes, or suturing.
  \item **Post-Operative Infection:** Spreading infection into the deep fascial spaces (cellulitis) or bone (osteomyelitis). Presents with severe pain, swelling, purulent discharge, trismus, and fever.
  \item **Nerve Injury:** Damage to the inferior alveolar nerve (IAN) or lingual nerve, primarily during mandibular third molar extractions. IAN damage causes temporary or permanent paresthesia, dysesthesia, or anesthesia of the lower lip, chin, and teeth on the affected side. Lingual nerve damage causes sensory loss or altered taste in the anterior two-thirds of the tongue.
  \item **Oroantral Communication (OAC):** Breach of the thin bone dividing the maxillary molar roots and the maxillary sinus. Can lead to chronic maxillary sinusitis or an oroantral fistula. Small defects close spontaneously; larger ones require surgical closure.
  \item **Tooth or Root Displacement:** Displacement of root tips into the maxillary sinus, submandibular space, or pterygomandibular space.
  \item **Mandibular Fracture:** A rare complication caused by excessive force during the extraction of deeply impacted third molars in osteoporotic or atrophic mandibles.
  \item **Damage to Adjacent Restorations:** Fracture or displacement of fillings/crowns on adjacent teeth.
\end{itemize}

\section*{Prognosis and Outcomes}
The prognosis for healing after a tooth extraction is excellent in healthy individuals, with the socket undergoing a predictable sequence of tissue repair:

Within the first 24 hours, a stable blood clot forms in the socket, serving as a scaffold for cellular migration. Over the next 2 to 7 days, granulation tissue replaces the clot, and epithelial cells migrate from the wound margins to cover the surface. By the third to fourth week, the socket is filled with rich connective tissue, and osteoid (unmineralized bone matrix) begins to deposit at the base and walls. Complete bone mineralization and remodelling take 3 to 6 months, transforming the osteoid into mature trabecular bone.

However, the extraction of a tooth inevitably leads to physiological resorption of the surrounding alveolar ridge. In the first year post-extraction, the alveolar ridge can lose up to 50\% of its width and height, with the majority of the loss occurring in the first three months. If socket preservation is not performed, this bone loss can compromise the aesthetic and functional outcome of future dental implants or prostheses.

Factors that can adversely affect prognosis and delay healing include advanced age, poor nutritional status (specifically Vitamin C and D deficiencies), systemic conditions (e.g., uncontrolled diabetes mellitus, immunosuppression), tobacco use (nicotine causes vasoconstriction, reducing blood supply and increasing the risk of dry socket), and non-compliance with post-operative instructions.

\section*{Prevention and Self-Care}
Adherence to post-operative self-care instructions is crucial for minimizing complications and ensuring rapid healing of the extraction socket:

\begin{itemize}
  \item **Immediate Post-Operative Care (First 24 Hours):** Bite firmly on the sterile gauze pack placed over the socket for 30--45 minutes. Avoid spitting, rinsing, sucking through a straw, or hot liquids for the first 24 hours to prevent dislodging the blood clot.
  \item **Swelling Control:** Apply cold compresses to the cheek (20 minutes on, 20 minutes off) for the first 24--48 hours.
  \item **Oral Hygiene:** Do not brush the surgical area on the first day. After 24 hours, rinse gently with warm salt water (1/2 teaspoon of salt in a glass of warm water) 4--5 times daily, particularly after meals, for 7--10 days.
  \item **Nutrition:** Eat soft, cool foods (yogurt, soup, pudding) and drink plenty of fluids. Avoid hard, crunchy, or hot foods.
  \item **Avoid Toxic Substances:** Strictly avoid smoking or using tobacco products for at least 72 hours (ideally 1 week), as smoking dramatically increases dry socket incidence. Avoid alcohol as it impairs clotting and interacts with analgesics.
  \item **Physical Activity:** Limit physical activity and avoid heavy lifting or exercise for the first 48 hours to prevent elevations in blood pressure that could dislodge the clot.
  \item **Long-Term Oral Hygiene:** Regular brushing, flossing, professional cleanings, and fluoride application to prevent decay and periodontal disease, minimizing the future need for extractions.
\end{itemize}

\section*{Sources and Further Reading}
The following reliable health organizations and peer-reviewed reference guides provide authoritative information on tooth extraction, oral surgery guidelines, and post-operative care:

\begin{itemize}
  \item **World Health Organization (WHO) -- Oral Health Program:** Comprehensive global resource on dental hygiene, disease prevention, and oral healthcare access. URL: \url{https://www.who.int/health-topics/oral-health}
  \item **American Dental Association (ADA) -- MouthHealthy:** Patient-focused educational resource covering tooth extraction procedures, post-operative care, and alternatives. URL: \url{https://www.mouthhealthy.org}
  \item **British Dental Association (BDA) -- Patient Support and Oral Health:** Authoritative guidelines on dental surgical procedures, patient consent, and recovery. URL: \url{https://bda.org}
  \item **National Institute of Dental and Craniofacial Research (NIDCR):** Research-based resources on wisdom teeth extraction, dry socket pathogenesis, and bone healing. URL: \url{https://www.nidcr.nih.gov}
  \item **Mayo Clinic -- Tooth Extraction Information:** Clinical overview of the indications, procedure steps, risks, and recovery guidelines. URL: \url{https://www.mayoclinic.org/tests-procedures/wisdom-tooth-extraction/about/pac-20395268}
  \item **Australian Dental Association (ADA):** National resource providing dental guidelines, patient fact sheets on oral surgery, and post-operative management. URL: \url{https://www.ada.org.au}
  \item **Cochrane Oral Health Group:** Systematic reviews on the efficacy of antibiotic prophylaxis in tooth extractions and treatments for alveolar osteitis. URL: \url{https://oralhealth.cochrane.org}
\end{itemize}